%% toc: P15 157 157.5,157.6,157.7
% Atomic, absolutely continuous and continuous singular parts of a real law
% (Durrett, PTE author version 5, p. 13 and thm. A.4.7; Feller II, secs. V.2--V.3),
% and the null modifications that preserve a conditional-expectation version.
\chapter{Mixed Real Laws and Null Modifications of Versions}
\label{sec:p15-w02b-p15-002-mixed-laws}

A real probability law can carry point masses, a part with a Lebesgue
density, and a remainder that has neither. This chapter proves that every
Borel probability law on the line splits uniquely into these three parts,
reads the split off its distribution function, and computes it for two
mixed laws. It also settles which changes of a conditional expectation on
a null set produce another version, with the witness showing that a
change on an arbitrary null event of the full sigma-algebra can fail.

\emph{Inputs.} A Borel probability measure \(\mu\) on \(\mathbb R\) with
distribution function \(F\); Lebesgue measure \(\lambda\) on the Borel
sets (\cref{thm:p15-lebesgue-unit-interval}); for the second part, a
probability space \((\Omega,\mathcal F_0,P)\), a sub-sigma-algebra
\(\mathcal G\) and a version of a conditional expectation
(\cref{def:p15-general-conditional-expectation}).
\emph{Conclusions.} The unique decomposition
\(\mu=\mu_{\mathrm d}+\mu_{\mathrm{ac}}+\mu_{\mathrm{sc}}\) into a purely
atomic, an absolutely continuous and a continuous singular measure; the
matching decomposition of \(F\), its mixture form and the corresponding
splitting of integrals; and an exact criterion for null modifications of
conditional-expectation versions.

The terms follow two held sources. Durrett calls a distribution function
absolutely continuous when it has a density, singular when its measure is
singular with respect to Lebesgue measure, and discrete when its measure
gives probability one to a countable set
\cite[p.~13]{DurrettPTEAuthor2019}. Feller defines atoms and atomic
distributions \cite[sec.~V.2, p.~137]{Feller1971ProbabilityII}, proves
the split of a probability distribution into an atomic and a continuous
component, which he calls the Jordan decomposition theorem
\cite[sec.~V.2, eqs.~(2.4)--(2.5), p.~138]{Feller1971ProbabilityII}, and
combines it with the Lebesgue decomposition relative to a given measure to
write every probability distribution as a mixture of an atomic, an
absolutely continuous and a continuous singular distribution
\cite[sec.~V.3, eq.~(3.9), p.~142]{Feller1971ProbabilityII}. Feller's
name refers to a different decomposition from the Jordan components of a
signed measure in \cref{thm:p15-jordan-decomposition}. The uniqueness
statements, the distribution-function form and the integration formula
below are proved here; Feller's text at the cited place states the
mixture without a uniqueness clause.

\section{Atoms, densities and the three parts of a real law}
\label{sec:p15-w02b-p15-002-three-parts}

\begin{hbdefinition}{def:p15-w02b-p15-002-law-types}{Atoms and the types of a finite Borel measure on the line}
Let \(\nu\) be a finite Borel measure on \(\mathbb R\). A point \(x\) is
an \emph{atom} of \(\nu\) if \(\nu(\{x\})>0\); write \(A_\nu\) for the set
of atoms. Say that \(\nu\) is
\begin{enumerate}
\item \emph{purely atomic} if \(\nu(\mathbb R\setminus S)=0\) for some
      countable set \(S\subseteq\mathbb R\);
\item \emph{atomless} if \(\nu(\{x\})=0\) for every real \(x\);
\item \emph{absolutely continuous} if \(\nu\ll\lambda\)
      (\cref{def:p15-absolute-continuity});
\item \emph{singular} if \(\nu\perp\lambda\)
      (\cref{def:p15-singularity});
\item \emph{continuous singular} if it is both atomless and singular.
\end{enumerate}
A countable set is a countable union of closed singletons, hence Borel.
The zero measure has every one of these types. For a probability measure,
purely atomic is Durrett's \emph{discrete}, atomless is Feller's
\emph{continuous}, and singular is Durrett's \emph{singular}.
Every purely atomic measure is singular in this sense, since a countable
set is \(\lambda\)-null; continuous singular measures are the singular
ones without atoms.
\index{atom!of a measure}\index{distribution!mixed}
\index{distribution!continuous singular}
\end{hbdefinition}

\begin{hblemma}{thm:p15-w02b-p15-002-atomic-split}{proved=annex:p15-w02b-p15-002-atomic-split}{Removing the atoms of a finite measure}
Let \(\nu\) be a finite Borel measure on \(\mathbb R\).
\begin{enumerate}
\item For each integer \(k\geq1\), at most \(k\,\nu(\mathbb R)\) points
      satisfy \(\nu(\{x\})\geq1/k\). Hence \(A_\nu\) is countable.
\item The formulas
      \[
       \nu_{\mathrm d}(E)\defeq\nu(E\cap A_\nu)
                       =\sum_{a\in E\cap A_\nu}\nu(\{a\}),
       \qquad
       \nu_{\mathrm c}(E)\defeq\nu(E\setminus A_\nu)
      \]
      define a purely atomic measure \(\nu_{\mathrm d}\) and an atomless
      measure \(\nu_{\mathrm c}\) with \(\nu=\nu_{\mathrm d}+\nu_{\mathrm c}\).
\item If \(\nu=\nu_1+\nu_2\) with Borel measures \(\nu_1\) purely atomic
      and \(\nu_2\) atomless, then \(\nu_1=\nu_{\mathrm d}\) and
      \(\nu_2=\nu_{\mathrm c}\).
\item \(\nu\) is purely atomic exactly when
      \(\nu(\mathbb R\setminus A_\nu)=0\), that is, when
      \(\nu_{\mathrm c}=0\).
\end{enumerate}
\end{hblemma}

\begin{hbtheorem}{thm:p15-w02b-p15-002-three-part}{proved=annex:p15-w02b-p15-002-three-part}{Atomic, absolutely continuous and continuous singular parts of a real law}
Let \(\mu\) be a Borel probability measure on \(\mathbb R\), let
\(F(x)=\mu((-\infty,x])\), and let \(A\) be the set of discontinuity
points of \(F\).
\begin{enumerate}
\item \emph{Existence.} The set \(A\) equals \(A_\mu\) and is
      countable, and there are Borel measures
      \(\mu_{\mathrm d},\mu_{\mathrm{ac}},\mu_{\mathrm{sc}}\) with
      \[
       \mu=\mu_{\mathrm d}+\mu_{\mathrm{ac}}+\mu_{\mathrm{sc}},
      \]
      \(\mu_{\mathrm d}\) purely atomic, \(\mu_{\mathrm{ac}}\)
      absolutely continuous and \(\mu_{\mathrm{sc}}\) continuous singular.
      Explicitly,
      \[
       \mu_{\mathrm d}(E)=\sum_{a\in E\cap A}\bigl(F(a)-F(a-)\bigr),
       \qquad
       \mu_{\mathrm{ac}}(E)=\int_E f\,d\lambda,
      \]
      for a real Borel \(f\geq0\) with \(\int f\,d\lambda\leq1\).
\item \emph{Uniqueness.} If \(\mu=\nu_1+\nu_2+\nu_3\) with Borel measures
      \(\nu_1\) purely atomic, \(\nu_2\) absolutely continuous and
      \(\nu_3\) continuous singular, then \(\nu_1=\mu_{\mathrm d}\),
      \(\nu_2=\mu_{\mathrm{ac}}\) and \(\nu_3=\mu_{\mathrm{sc}}\).
      Any two Borel densities of \(\mu_{\mathrm{ac}}\) with respect to
      \(\lambda\) agree \(\lambda\)-almost everywhere.
\item \emph{Distribution functions.} Put
      \[
       \begin{aligned}
       F_{\mathrm d}(x)&=\sum_{a\in A,\ a\leq x}\bigl(F(a)-F(a-)\bigr),\\
       F_{\mathrm{ac}}(x)&=\int_{(-\infty,x]}f\,d\lambda,\\
       F_{\mathrm{sc}}(x)&=\mu_{\mathrm{sc}}((-\infty,x]).
       \end{aligned}
      \]
      Then \(F=F_{\mathrm d}+F_{\mathrm{ac}}+F_{\mathrm{sc}}\). Each of the
      three is bounded, nondecreasing and right continuous with limit zero
      at \(-\infty\); \(F_{\mathrm{ac}}\) and \(F_{\mathrm{sc}}\) are
      continuous; and \(F_{\mathrm d}(x)-F_{\mathrm d}(x-)=F(x)-F(x-)\) for
      every \(x\). If \(F=G_1+G_2+G_3\), where each \(G_i\) is bounded,
      nondecreasing and right continuous with limit zero at \(-\infty\),
      and the Stieltjes measures \(\mu_{G_1},\mu_{G_2},\mu_{G_3}\) of
      \cref{thm:p15-stieltjes-construction} are respectively purely
      atomic, absolutely continuous and continuous singular, then
      \(G_1=F_{\mathrm d}\), \(G_2=F_{\mathrm{ac}}\) and
      \(G_3=F_{\mathrm{sc}}\).
\item \emph{Mixture form.} The weights
      \(p_{\mathrm d}=\mu_{\mathrm d}(\mathbb R)\),
      \(p_{\mathrm{ac}}=\mu_{\mathrm{ac}}(\mathbb R)\) and
      \(p_{\mathrm{sc}}=\mu_{\mathrm{sc}}(\mathbb R)\) are nonnegative
      with sum one. For each positive weight \(p_\bullet\), the measure
      \(\mu_\bullet/p_\bullet\) is a Borel probability law of the same
      type, with distribution function \(F_\bullet/p_\bullet\), and
      \(F\) is the sum of \(p_\bullet\,(F_\bullet/p_\bullet)\) over the
      positive weights. A part with zero weight is the zero measure.
\item \emph{Integration.} For every Borel \(g:\mathbb R\to[0,\infty]\),
      \[
       \int g\,d\mu=\sum_{a\in A}g(a)\bigl(F(a)-F(a-)\bigr)
                    +\int gf\,d\lambda+\int g\,d\mu_{\mathrm{sc}}.
      \]
      If \(g\) is real Borel with \(\int|g|\,d\mu<\infty\), the three
      terms are finite, the series converges absolutely, and the same
      identity holds.
\end{enumerate}
\end{hbtheorem}

In Feller's notation the normalized atomic part is his atomic component
\(F_a\) of \cite[sec.~V.2, eq.~(2.4), p.~138]{Feller1971ProbabilityII},
with \(p_{\mathrm d}\) his total atomic weight \(p\); his continuous
component is \((F_{\mathrm{ac}}+F_{\mathrm{sc}})/(1-p_{\mathrm d})\) when
\(p_{\mathrm d}<1\).

\begin{hbcorollary}{thm:p15-w02b-p15-002-regimes}{proved=annex:p15-w02b-p15-002-regimes}{Four regimes of a real law}
For a Borel probability law \(\mu\) on \(\mathbb R\) with the parts of
\cref{thm:p15-w02b-p15-002-three-part}:
\begin{enumerate}
\item \(\mu\) is discrete exactly when \(\mu=\mu_{\mathrm d}\), exactly
      when \(\sum_{a\in A}(F(a)-F(a-))=1\);
\item \(F\) is continuous exactly when \(\mu_{\mathrm d}=0\);
\item \(\mu\) has a Lebesgue density exactly when \(\mu=\mu_{\mathrm{ac}}\);
\item \(\mu\) is singular with respect to \(\lambda\) exactly when
      \(\mu_{\mathrm{ac}}=0\);
\item the fair-digit Cantor law \(\eta\) of
      \cref{def:p15-cantor-law-construction} satisfies
      \(\eta=\eta_{\mathrm{sc}}\).
\end{enumerate}
A law is \emph{mixed} when at least two of its weights are positive; such
a law lies in none of the classes of items~1, 3 and~5, namely the laws with
\(\mu=\mu_{\mathrm d}\), with \(\mu=\mu_{\mathrm{ac}}\), and with
\(\mu=\mu_{\mathrm{sc}}\), the class to which item~5 assigns \(\eta\).
\end{hbcorollary}

\begin{hbmethod}{def:p15-w02b-p15-002-procedure}{Identifying the three parts}
\begin{enumerate}
\item List the jumps of \(F\). Each discontinuity point \(a\) contributes
      the atom \(F(a)-F(a-)\) (\cref{thm:p15-cdf-laws}); their sum is
      \(p_{\mathrm d}\) and their cumulative sum is \(F_{\mathrm d}\).
\item Subtract: \(F-F_{\mathrm d}=F_{\mathrm{ac}}+F_{\mathrm{sc}}\) is
      continuous and nondecreasing.
\item Propose a Borel \(f\geq0\) and an atomless measure
      \(\sigma\) concentrated on a \(\lambda\)-null Borel set with
      \(F-F_{\mathrm d}=\int_{(-\infty,\cdot\,]}f\,d\lambda
      +\sigma((-\infty,\cdot\,])\).
\item Verify the three types. Uniqueness in
      \cref{thm:p15-w02b-p15-002-three-part} then identifies
      \(\mu_{\mathrm{ac}}\) and \(\mu_{\mathrm{sc}}=\sigma\); nothing
      further needs checking.
\end{enumerate}
Side conditions: the decomposition is relative to Lebesgue measure on
the Borel sets of \(\mathbb R\), \(f\) is determined only
\(\lambda\)-almost everywhere, and step~3 is a verification step. When the
law is given as a sum of a point-mass part, a density part and a
singular atomless part, step~4 applies directly.
\end{hbmethod}

\begin{hbtheorem}{thm:p15-w02b-p15-002-mixed-examples}{proved=annex:p15-w02b-p15-002-mixed-examples}{A censored uniform variable and a law with all three parts}
\begin{enumerate}
\item On the uniform probability space of
      \cref{thm:p15-lebesgue-unit-interval}, let \(X=\min(U,1/2)\).
      Its distribution function is
      \[
       F_X(x)=
       \begin{cases}
       0,&x<0,\\
       x,&0\leq x<1/2,\\
       1,&x\geq1/2.
       \end{cases}
      \]
      It has one jump, of size \(1/2\) at \(x=1/2\). The parts of its law
      are \(\mu_{\mathrm d}=\tfrac12\delta_{1/2}\),
      \(\mu_{\mathrm{ac}}(E)=\lambda(E\cap(0,1/2))\) with density
      \(\mathbf1_{(0,1/2)}\), and \(\mu_{\mathrm{sc}}=0\). In mixture form,
      \(p_{\mathrm d}=p_{\mathrm{ac}}=1/2\): the law is the even mixture of
      the point mass at \(1/2\) and the uniform law on \((0,1/2)\) with
      density \(2\cdot\mathbf1_{(0,1/2)}\). Moreover
      \[
       \mathbb E[X]=\tfrac12\cdot\tfrac12+\int_{(0,1/2)}x\,d\lambda(x)
                   =\tfrac14+\tfrac18=\tfrac38,
      \]
      in agreement with the direct computation of
      \(\mathbb E[\min(U,1/2)]\).
\item Let \(\eta\) be the Cantor law of
      \cref{thm:p15-cantor-singular-law}, with distribution function
      \(F_\eta\), and let
      \[
       \rho=\tfrac14\delta_2+\tfrac12\lambda(\,\cdot\,\cap(0,1))
             +\tfrac14\eta .
      \]
      Then \(\rho\) is a Borel probability law with
      \(\rho_{\mathrm d}=\tfrac14\delta_2\),
      \(\rho_{\mathrm{ac}}=\tfrac12\lambda(\,\cdot\,\cap(0,1))\) and
      \(\rho_{\mathrm{sc}}=\tfrac14\eta\). Its distribution function
      \[
       F_\rho(x)=\tfrac14\mathbf1_{[2,\infty)}(x)
                 +\tfrac12\min(\max(x,0),1)+\tfrac14F_\eta(x)
      \]
      has exactly one jump, of size \(1/4\) at \(x=2\); it equals
      \(x/2+1/8\) on \([1/3,2/3]\); and \(F_\rho(1)-F_\rho(0)=3/4\),
      whereas the density part contributes only \(1/2\) on \([0,1]\).
      The Cantor set \(C\) has \(\lambda(C)=0\) and \(\rho(C)=1/4\).
\end{enumerate}
\end{hbtheorem}

\begin{hbexample}{def:p15-w02b-p15-002-reading}{Reading the two distribution functions}
For the censored variable, the graph of \(F_X\) rises with slope one on
\([0,1/2)\), reaching the left limit \(F_X(1/2-)=1/2\), and then jumps to
one. The jump gives the atomic part and the slope gives the density part;
nothing is left over, so the singular part vanishes. Both weights are
\(1/2\), and the mean \(3/8\) is the weighted average of the point-mass
mean \(1/2\) and the uniform mean \(1/4\).

For \(\rho\), the only jump is at \(2\). On the removed middle third
\((1/3,2/3)\) the Cantor distribution function is constant, so
\(F_\rho\) increases there exactly as the density part does, with slope
\(1/2\); for instance \(F_\rho(1/2)=3/8\). Across \([0,1]\), however,
\(F_\rho\) increases by \(3/4\) without any jump, while the density part
accounts for \(1/2\). The missing \(1/4\) is carried by the Cantor set,
which has Lebesgue measure zero: it is the continuous singular part, and
no inspection of jumps alone detects it.
\end{hbexample}

\begin{hbwarning}{def:p15-w02b-p15-002-limits}{What the decomposition depends on and what it does not give}
\begin{enumerate}
\item The split between the absolutely continuous and the singular parts
      is taken relative to Lebesgue measure. Relative to another base
      measure it can change, as in item~1 of
      \cref{thm:p15-radon-nikodym-examples}, where adding a point mass to
      the base makes the whole measure absolutely continuous.
\item Terminology differs between the sources. Durrett's singular laws
      include the discrete ones \cite[p.~13]{DurrettPTEAuthor2019};
      Feller's continuous distributions are those without atoms, a
      condition on the distribution function and not on any density
      \cite[sec.~V.2, p.~138]{Feller1971ProbabilityII}. Only the
      continuous singular part excludes both atoms and densities.
\item The jump set can be dense. Durrett's Example~1.2.9 places positive
      masses at all rationals \cite[p.~14]{DurrettPTEAuthor2019}, and
      Feller gives an atomic distribution for which every point of the
      closed unit interval is a point of increase
      \cite[sec.~V.2, p.~137]{Feller1971ProbabilityII}. Then
      \(F_{\mathrm d}\) is discontinuous on a dense set and the parts
      cannot be read off from finitely many points.
\item The density \(f\) is determined only \(\lambda\)-almost everywhere;
      \(F_{\mathrm{ac}}\) does not single out a pointwise representative.
\item No description of \(F_{\mathrm{sc}}\) or \(F_{\mathrm{ac}}\) by
      derivatives of \(F\) is asserted here; the procedure above verifies a
      proposed split rather than extracting \(f\) from \(F\) by
      differentiation.
\item The statement concerns Borel measures on \(\mathbb R\). The atomic
      split uses that singletons are measurable and that the measure is
      finite: a nonzero multiple of counting measure on an uncountable set
      has uncountably many atoms.
\end{enumerate}
\end{hbwarning}

\section{Null events and versions of conditional expectations}
\label{sec:p15-w02b-p15-002-null-versions}

\Cref{thm:p15-conditional-limits} assembles its constructions by changing
versions on null sets that belong to the conditioning sigma-algebra, and
asserts that arbitrary changes on a null event of the full
sigma-algebra can destroy measurability with respect to the conditioning
sigma-algebra. The next theorem gives the exact criterion and a witness.

\begin{hbtheorem}{thm:p15-w02b-p15-002-null-modification}{proved=annex:p15-w02b-p15-002-null-modification}{Null modifications of a version}
Let \((\Omega,\mathcal F_0,P)\) be a probability space,
\(\mathcal G\subseteq\mathcal F_0\) a sub-sigma-algebra, \(X\) an
integrable real random variable, and \(Z\) a version of
\(E(X\mid\mathcal G)\).
\begin{enumerate}
\item A real \(\mathcal F_0\)-measurable \(Z'\) with \(Z'=Z\) almost
      surely is a version of \(E(X\mid\mathcal G)\) exactly when \(Z'\) is
      \(\mathcal G\)-measurable.
\item If \(N\in\mathcal G\), \(P(N)=0\), and \(W\) is real and
      \(\mathcal G\)-measurable, then
      \(Z\mathbf1_{N^c}+W\mathbf1_N\) is a version.
\item If \(N\in\mathcal F_0\) and \(P(N)=0\), then
      \(Z+\mathbf1_N\) is a version exactly when \(N\in\mathcal G\).
\item If \(\mathcal G\) contains every \(P\)-null set of
      \(\mathcal F_0\), then every real \(\mathcal F_0\)-measurable
      \(Z'\) with \(Z'=Z\) almost surely is a version.
\item Let \(\mathcal G=\{\varnothing,\Omega\}\). The constant
      \(\mathbb E[X]\) is the only version, pointwise, and \(\varnothing\)
      is the only \(\mathcal G\)-measurable null set. If
      \(N\in\mathcal F_0\) is nonempty with \(P(N)=0\), then
      \(\mathbb E[X]+\mathbf1_N\) equals the version outside a null event
      but is not a version. On the uniform probability space with
      \(X=U\) and \(N=\{1/2\}\), the version is the constant \(1/2\), and
      the function equal to \(1/2\) except for the value \(3/2\) at
      \(u=1/2\) is not a version.
\item On the uniform probability space let
      \(L=\{U\leq1/2\}=(0,1/2]\) and
      \(\mathcal G=\{\varnothing,\Omega_U,L,L^c\}\). Then
      \[
       E(U\mid\mathcal G)=\tfrac14\mathbf1_L+\tfrac34\mathbf1_{L^c}
      \]
      is the only version, pointwise. Changing it at the single point
      \(u=1/4\) gives a function equal to it almost surely that is not a
      version.
\end{enumerate}
\end{hbtheorem}

Items~1 and~2 are parts~(ii) and~(i) of
\cref{thm:p15-w02b-p15-008-null-modification}, item~1 in the case of an
\(\mathcal F_0\)-measurable modification, and item~5 includes, as the case
of Lebesgue measure on \([0,1]\), that theorem's part~(iii). Item~2 is the
operation used in the proof of \cref{thm:p15-conditional-limits}; items~3 and~5 show that the
restriction to null sets in \(\mathcal G\) is needed, and item~4 shows
that it can be dropped exactly when that restriction is vacuous because
\(\mathcal G\) already contains all null events. Items~5 and~6 also show
that when every nonempty set of \(\mathcal G\) has positive probability,
almost-everywhere uniqueness of the version is pointwise uniqueness.

\begin{hbwarning}{def:p15-w02b-p15-002-version-limits}{Almost-sure equality is not the defining property of a version}
A version is fixed by two requirements: measurability with respect to
\(\mathcal G\) and the integral identity on every event of \(\mathcal G\).
Almost-sure equality to a version preserves the second requirement but not
the first. In a statement that holds outside a null event, the exceptional
event must therefore be chosen in \(\mathcal G\) whenever the modified
object is to remain a version; \cref{thm:p15-conditional-existence} needs
no completeness of \(\mathcal G\) for exactly this reason. Item~4 of
\cref{thm:p15-w02b-p15-002-null-modification} is the only general
condition under which the distinction disappears, and it alters
\(\mathcal G\).
\end{hbwarning}

\section{Relationships}
\label{sec:p15-w02b-p15-002-relationships}

\Cref{thm:p15-w02b-p15-002-three-part} combines three earlier results:
the countability of jumps and the point-mass formula of
\cref{thm:p15-cdf-laws}, the two-part Lebesgue decomposition of
\cref{thm:p15-lebesgue-decomposition}, applied to the atomless remainder
with Lebesgue measure as base, and the density representation of
\cref{thm:p15-radon-nikodym}. The Cantor law of
\cref{thm:p15-cantor-singular-law} is the case of a pure continuous
singular part. Every law in the theorem, including each mixture, is
realized as the law of the quantile variable \(Q_F\) on the uniform
probability space (\cref{thm:p15-cdf-existence}); in the censored example
that variable has the same law as \(\min(U,1/2)\). Countable atomic models
(\cref{thm:p15-countable-atomic}) carry exactly the purely atomic laws
when their outcomes are real.

The same distinction between a representation and its minimal or
canonical form appears for signed measures on the line: the cumulative
function of a finite signed measure can be written as \(F-G\) with
Stieltjes measures that exceed the Jordan components by a common finite
measure, and the choice \(F=G=\mathbf1_{[0,\infty)}\) represents the
zero signed measure with \(\mu_F=\mu_G=\delta_0\). This witness and the
exact description are \cref{thm:p15-w02b-p15-005-stieltjes-excess}.

For conditional laws, a regular conditional distribution defined only
outside a null event shows the failure of
\cref{thm:p15-w02b-p15-002-null-modification} in a kernel setting: filling
in the exceptional event by a fixed value can destroy measurability
(\cref{thm:p15-w02b-p15-008-ae-counterexample}, item~4).

\emph{Sources.} Definitions of absolutely continuous, singular and
discrete distribution functions, the Cantor example and the dense-jump
example: Durrett, author version~5, dated 11 January 2019
\cite[pp.~13--14]{DurrettPTEAuthor2019}. Lebesgue decomposition for
sigma-finite measures: \cite[thm.~A.4.7, p.~469]{DurrettPTEAuthor2019}.
Atoms, the atomic component and the atomic and continuous split:
\cite[sec.~V.2, pp.~137--138]{Feller1971ProbabilityII}. The three-part
mixture: \cite[sec.~V.3, p.~142]{Feller1971ProbabilityII}. Definition of
conditional expectation and its almost-everywhere uniqueness:
\cite[pp.~205--206]{DurrettPTEAuthor2019}.
